\documentclass[11pt,a4paper]{article}
\usepackage[margin=1in]{geometry}
\usepackage{amsmath,amssymb,amsthm}
\usepackage{graphicx}
\usepackage{microtype}
\usepackage[hidelinks]{hyperref}

\newtheorem{theorem}{Theorem}
\newtheorem{lemma}[theorem]{Lemma}
\newtheorem{remark}[theorem]{Remark}
\newcommand{\dist}{\operatorname{dist}}

\title{Generalized lushness passes to all\\Bochner $L_\infty$ spaces}
\author{A full answer to one open direction in arXiv:2007.10589}
\date{17 August 2026}

\begin{document}
\maketitle

\noindent\textbf{Status.}
Candidate full solution, likely valid, pending human review.  If a real Banach
space $X$ is generalized-lush, then $L_\infty(\mu,X)$ is generalized-lush for
every measure space.  This resolves the previously unknown forward
$L_\infty$ implication in the source.  The three converse questions stated
beside it are not claimed here.

\section{Source question}

For $x^*\in S_{X^*}$ and $0<\alpha<1$, write
\[
 S_X(x^*,\alpha)=\{x\in B_X:x^*(x)>1-\alpha\}.
\]
A Banach space $X$ is \emph{generalized-lush} (GL) if for every
$x\in S_X$ and $\alpha>0$ there is $x^*\in S_{X^*}$ such that
$x\in S_X(x^*,\alpha)$ and
\[
 \dist(y,S_X(x^*,\alpha))+
 \dist(y,-S_X(x^*,\alpha))<2+\alpha
 \qquad(y\in S_X).
\]

Zhao and Tan record that GL is known to pass from $X$ to $C(K,X)$ and
$L_1(\mu,X)$, but explicitly state that the corresponding assertion for
$L_\infty(\mu,X)$ was unknown~\cite{ZT}.

\begin{center}
\includegraphics[width=\textwidth]{figures/open_question_crop.png}
\end{center}

\section{Proof intuition}

Fix $f\in S_{L_\infty(\mu,X)}$.  Choose a Boolean ultrafilter concentrating
on a set where $\|f(t)\|$ is almost one.  The normalized values of $f$ have
separable range, so approximate them uniformly by a countably valued map.
At each of those countably many values choose a GL slice functional on $X$.
Taking the ultrafilter limit of the resulting scalar pairings defines one
functional on $L_\infty(\mu,X)$ whose slice contains $f$.

For an arbitrary $g$, countable approximation again makes pointwise GL
witnesses measurable.  Their two pointwise distances have sum below
$2+o(1)$, but taking two separate essential suprema could destroy that
bound.  The key step is to divide the first distance into finitely many
short intervals.  The ultrafilter selects one interval.  Patch the witnesses
only on that selected region and leave $g$ unchanged elsewhere.  On the
selected region the first distance has almost no oscillation, so the two
global suprema still add to less than $2+\varepsilon$.

\section{Two elementary lemmas}

\begin{lemma}[The GL estimate holds on the ball]\label{lem:ball}
Let $S$ be a slice of $B_X$.  If
\[
 D(y):=\dist(y,S)+\dist(y,-S)<c\qquad(y\in S_X),
\]
then $D(y)<c$ for every $y\in B_X$.
\end{lemma}

\begin{proof}
The function $D$ is convex and even.  If $y=rz$ with $z\in S_X$ and
$0\le r\le1$, then
\[
 y=\frac{1+r}{2}z+\frac{1-r}{2}(-z).
\]
Convexity and evenness give $D(y)\le D(z)<c$.
\end{proof}

Let $\mathfrak B=\Sigma/\mathcal N_\mu$ be the measure algebra.  If
$\mathcal U$ is an ultrafilter on $\mathfrak B$ and $q$ is a bounded real
measurable function, $\lim_{\mathcal U}q$ denotes its ultrafilter limit.
It is unchanged by null modifications and preserves sums and scalar
multiples.

\begin{lemma}[A countably valued field gives a functional]
\label{lem:functional}
Suppose $\phi:\Omega\to S_{X^*}$ is constant on each member of a countable
measurable partition.  Then
\[
 F_0(h):=\lim_{\mathcal U}\phi(t)(h(t))
 \qquad(h\in L_\infty(\mu,X))
\]
is a well-defined linear functional of norm at most one.  If
$[E]\in\mathcal U$ and $\phi(t)(h(t))\ge c$ almost everywhere on $E$, then
$F_0(h)\ge c$.  The analogous upper estimate also holds.
\end{lemma}

\begin{proof}
For strongly measurable $h$, the scalar function
$t\mapsto\phi(t)(h(t))$ is measurable because $\phi$ is countably valued.
It is bounded in absolute value by $\|h\|_\infty$.  Linearity and the norm
bound follow from the corresponding properties of ultrafilter limits.
The last assertion follows because an ultrafilter limit ignores the
complement of every member of $\mathcal U$.
\end{proof}

\section{The full \texorpdfstring{$L_\infty$}{L-infinity} stability theorem}

\begin{theorem}\label{thm:main}
Let $(\Omega,\Sigma,\mu)$ be a nonzero measure space and let $X$ be a real
Banach space.  If $X$ is generalized-lush, then the Bochner space
$L_\infty(\mu,X)$ is generalized-lush.
\end{theorem}

\begin{proof}
Fix $f\in S_{L_\infty(\mu,X)}$ and $0<\varepsilon<1$, and put
$\delta=\varepsilon/20$.  The measurable set
\[
 A=\{t:\|f(t)\|>1-\delta\}
\]
is non-null.  Extend the filter generated by $[A]$ to an ultrafilter
$\mathcal U$ on $\mathfrak B$.

On $A$, the map $u(t)=f(t)/\|f(t)\|$ is strongly measurable and has
separable range.  Hence there are a measurable partition
$A=\bigcup_{n\ge1}A_n$ and points $x_n\in S_X$ such that
\begin{equation}\label{eq:approx-f}
 \|u(t)-x_n\|<\delta\qquad(t\in A_n).
\end{equation}
For every $n$, generalized lushness supplies
$\phi_n\in S_{X^*}$ such that
\begin{equation}\label{eq:phi-n}
 \phi_n(x_n)>1-\delta
\end{equation}
and, with $S_n=S_X(\phi_n,\delta)$,
\begin{equation}\label{eq:GL-n}
 \dist(y,S_n)+\dist(y,-S_n)<2+\delta
 \qquad(y\in B_X),
\end{equation}
where Lemma~\ref{lem:ball} was used for the ball.

Define $\phi(t)=\phi_n$ on $A_n$ and give it any fixed value in $S_{X^*}$
off $A$.  Lemma~\ref{lem:functional} defines $F_0\in
L_\infty(\mu,X)^*$ with $\|F_0\|\le1$.  On $A_n$,
\[
 \phi(t)(f(t))
 =\|f(t)\|\phi_n(u(t))
 >(1-\delta)(1-2\delta)>1-3\delta
\]
by \eqref{eq:approx-f} and \eqref{eq:phi-n}.  Since $[A]\in\mathcal U$,
\begin{equation}\label{eq:Ff}
 F_0(f)\ge1-3\delta>0.
\end{equation}
Thus $c:=\|F_0\|>0$; put $F=F_0/c\in S_{L_\infty(\mu,X)^*}$.
Because $c\le1$, \eqref{eq:Ff} gives
\[
 F(f)>1-\varepsilon.
\]

It remains to verify the distance condition for an arbitrary
$g\in S_{L_\infty(\mu,X)}$.  Strong measurability gives a countably valued
measurable map $s(t)=y_m$ on a partition $(C_m)_{m\ge1}$, with
$y_m\in B_X$ and
\begin{equation}\label{eq:approx-g}
 \|g(t)-s(t)\|<\delta\quad\text{almost everywhere}.
\end{equation}
For each pair $(n,m)$, \eqref{eq:GL-n} permits choices
$a_{nm}\in S_n$ and $b_{nm}\in-S_n$ satisfying
\begin{equation}\label{eq:point-witness}
 \|y_m-a_{nm}\|+\|y_m-b_{nm}\|<2+2\delta.
\end{equation}
On $A_n\cap C_m$ set $a(t)=a_{nm}$ and $b(t)=b_{nm}$.  These maps are
strongly measurable, take values in $B_X$, and satisfy
\begin{align}
 \phi(t)(a(t))&>1-\delta, &
 \phi(t)(b(t))&<-1+\delta,\label{eq:slices-pointwise}\\
 d_+(t)+d_-(t)&<2+4\delta,
 &d_\pm(t)&:=\|g(t)-a(t)\|,\ \|g(t)-b(t)\|,
 \label{eq:distance-pointwise}
\end{align}
on $A$, by \eqref{eq:approx-g} and \eqref{eq:point-witness}.

Partition $A$ into finitely many measurable sets according to which interval
of length $\delta$ contains $d_+(t)$; the last interval is truncated at two.
Exactly one member $E$ of this finite partition has $[E]\in\mathcal U$.
For some integer $j\ge0$ we therefore have on $E$
\[
 j\delta\le d_+(t)<(j+1)\delta,
 \qquad
 d_-(t)<2+4\delta-j\delta.
\]
Define
\[
 h_+(t)=\begin{cases}a(t),&t\in E,\\g(t),&t\notin E,
 \end{cases}
 \qquad
 h_-(t)=\begin{cases}b(t),&t\in E,\\g(t),&t\notin E.
 \end{cases}
\]
Both functions lie in the unit ball.  Since $[E]\in\mathcal U$,
Lemma~\ref{lem:functional} and \eqref{eq:slices-pointwise} give
\[
 F_0(h_+)\ge1-\delta,
 \qquad F_0(h_-)\le-1+\delta.
\]
Dividing by $c\le1$ only improves these signed inequalities, so
\[
 h_+\in S_{L_\infty}(F,\varepsilon),
 \qquad h_-\in-S_{L_\infty}(F,\varepsilon).
\]
Finally, the patching and the selected distance bin give
\begin{align*}
 \|g-h_+\|_\infty+
 \|g-h_-\|_\infty
 &\le (j+1)\delta+\bigl(2+4\delta-j\delta\bigr)\\
 &=2+5\delta<2+\varepsilon.
\end{align*}
The same functional $F$ works for every $g$; hence
$L_\infty(\mu,X)$ is GL.
\end{proof}

\section{Upgrade attempts, novelty, and scope}

The first route used an ordinary average over a positive-measure set.  It
fails because the sum of two pointwise distances does not control the sum of
their two essential suprema.  A second route worked when $X$ was
finite-dimensional: an ultrafilter limit of every bounded $X$-valued
function then lies in the compact unit ball, so one can patch near that
limit.  The decisive upgrade removes compactness.  It uses a countably
valued field of slice functionals and lets the Boolean ultrafilter select a
\emph{finite distance bin}; this controls both suprema without requiring an
$X$-valued ultralimit.  It also shows that no $\sigma$-finiteness assumption
is needed.

A bounded primary-source search covered the exact source statement,
``generalized-lush'', ``GL-space'', ``$L_\infty(\mu,X)$'', and later arXiv
citations.  It found the source, the original paper proving stability under
$C(K,X)$ and $\ell_\infty$-sums, and work on property $(**)$, but no primary
paper proving this implication.  Novelty confidence is moderate; correctness
confidence is high.

The proof resolves only the forward implication from $X$ to
$L_\infty(\mu,X)$.  It does not show that GL descends from $C(K,X)$,
$L_1(\mu,X)$, or $L_\infty(\mu,X)$ to a nonseparable $X$.  Direct restriction
and averaging attempts at those converses lose the single slice functional
that must work simultaneously for the entire nonseparable sphere.  Those
questions remain outside the claim of this packet.

\begin{thebibliography}{1}
\bibitem{ZT}
K.~Zhao and D.~Tan,
\emph{A property in vector-valued function spaces},
arXiv:2007.10589.
\end{thebibliography}

\end{document}
